% Part: propositional-logic
% Chapter: syntax-and-semantics
% Section: introduction

\documentclass[../../../include/open-logic-section]{subfiles}

\begin{document}

\olfileid[hi]{pl}{syn}{int}

\olsection{परिचय}

प्रतिज्ञप्ति तर्कशास्त्र उन !!{formula}s का अध्ययन करता है जो
!!{propositional variable}s से प्रतिज्ञप्ति संयोजकों $\lnot$, $\land$,
$\lor$, $\lif$ और $\liff$ का उपयोग करके बनाए जाते हैं। सहज रूप से,
!!a{propositional variable}~$p$ किसी ऐसे वाक्य या प्रतिज्ञप्ति का स्थान लेता
है जो सत्य या असत्य हो। किसी !!{propositional variable} का ``सत्य-मूल्य''
किसी !!a{formula} में निर्धारित होते ही, उससे प्रतिज्ञप्ति संयोजकों द्वारा बने
सभी !!{formula}s का सत्य-मूल्य भी निर्धारित हो जाता है। हम कहते हैं कि
प्रतिज्ञप्ति तर्कशास्त्र \emph{सत्य-फलनीय} है, क्योंकि उसका अर्थविज्ञान
सत्य-मूल्यों के फलनों द्वारा दिया जाता है। विशेषतः प्रतिज्ञप्ति तर्कशास्त्र में
हम सत्यता और असत्यता के किसी अतिरिक्त निर्धारण पर विचार नहीं करते; जैसे, कोई
बात केवल सांयोगिक रूप से सत्य होने के बजाय अनिवार्य रूप से सत्य है या नहीं,
किसी बात का सत्य होना ज्ञात है या नहीं, अथवा कोई बात अभी सत्य है न कि पहले
सत्य थी या आगे सत्य होगी। हम केवल दो सत्य-मूल्य—सत्य ($\True$) और असत्य
($\False$)—मानते हैं; इसलिए हम इस संभावना को चर्चा से बाहर रखते हैं कि कोई
कथन न सत्य हो न असत्य, या केवल आधा सत्य हो। हम केवल ऐसे संयोजकों पर भी ध्यान
देते हैं जिनसे बने !!a{formula} का सत्य-मूल्य उसके अवयवों के सत्य-मूल्यों से
पूरी तरह निर्धारित होता है (न कि, मसलन, उसके अर्थ से)। विशेषतः, अंग्रेज़ी के
सशर्त वाक्यों का सत्य-मूल्य इस अर्थ में सत्य-फलनीय है या नहीं, यह विवादास्पद
है। शाब्दिक आपादन~$\lif$ सत्य-फलनीय है; अन्य तर्क-पद्धतियाँ ऐसे सशर्त वाक्यों
का अध्ययन करती हैं जो सत्य-फलनीय नहीं होते।

सत्य-फलनीय प्रतिज्ञप्ति तर्कशास्त्र के सिद्धांत और अधिसिद्धांत को विकसित करने
के लिए हमें पहले उसकी अभिव्यक्तियों के वाक्यविन्यास और अर्थविज्ञान को परिभाषित
करना होगा। हम !!{formula}s को !!{propositional variable}s से संयोजकों का
उपयोग करके बनाने का एक ढंग बताएँगे। वैकल्पिक परिभाषाएँ भी संभव हैं। अन्य
पद्धतियाँ भिन्न प्रतीक चुनेंगी, संयोजकों के भिन्न समुच्चयों को आदिम मानेंगी,
और कोष्ठकों का भिन्न प्रकार से उपयोग करेंगी (या तथाकथित पोलिश संकेतन की तरह
उनका बिल्कुल उपयोग नहीं करेंगी)। फिर भी सभी दृष्टिकोणों में यह बात समान है कि
निर्माण-नियम !!{formula}s के समुच्चय को \emph{आगमनात्मक रूप से} परिभाषित करते
हैं। परिभाषा ठीक ढंग से दी गई हो तो निर्माण-नियमों के अनुसार प्रत्येक
अभिव्यक्ति मूलतः केवल एक ही प्रकार से बन सकती है। अभिव्यक्तियों को
\emph{अद्वितीय रूप से पठनीय} बनाने वाली आगमनात्मक परिभाषा के कारण हम उसी
विधि—आगमनात्मक परिभाषा—से इन अभिव्यक्तियों को अर्थ दे सकते हैं।

अभिव्यक्तियों को अर्थ देना अर्थविज्ञान का क्षेत्र है। प्रतिज्ञप्ति तर्कशास्त्र
के अर्थविज्ञान की केंद्रीय धारणा किसी !!a{valuation} में संतुष्टि की धारणा है।
!!^a{valuation}~$\pAssign{v}$ !!{propositional variable}s को सत्य-मूल्य
$\True$, $\False$ देता है। कोई भी !!{valuation} सत्य-मूल्य $\pValue{v}(!A)$
निर्धारित करता है—किसी भी !!{formula}~$!A$ के लिए। कोई !!^a{formula},
!!a{valuation}~$\pAssign{v}$ में तभी और केवल तभी संतुष्ट होता है जब
$\pValue{v}(!A) = \True$—इसे हम $\pSat{v}{!A}$ लिखते हैं। इस संबंध को भी
तार्किक संयोजकों के सत्य-फलनों का उपयोग करते हुए~$!A$ की संरचना पर आगमन द्वारा
परिभाषित किया जा सकता है; उदाहरणार्थ, $!A \land !B$ की संतुष्टि को $!A$ और~$!B$
की संतुष्टि (या असंतुष्टि) के आधार पर परिभाषित किया जाता है।

वाक्यों के लिए संतुष्टि-संबंध $\pSat{v}{!A}$ के आधार पर हम पुनरुक्ति, तार्किक
परिणाम और संतुष्टता की मूल अर्थवैज्ञानिक धारणाएँ परिभाषित कर सकते हैं। कोई
!!^a{formula} पुनरुक्ति है, $\Entails !A$, यदि प्रत्येक !!{valuation} उसे
संतुष्ट करता है; अर्थात् $\pValue{v}(!A) = \True$ किसी भी~$\pAssign{v}$ के
लिए। !!{formula}s का कोई समुच्चय किसी सूत्र को तार्किकतः निष्पन्न करता है,
$\Gamma \Entails !A$, यदि प्रत्येक !!{valuation}, जो~$\Gamma$ के सभी
!!{formula}s को संतुष्ट करता है, $!A$ को भी संतुष्ट करता है। और
!!{formula}s का कोई समुच्चय संतुष्ट है यदि कोई !!{valuation} उसमें उपस्थित सभी
!!{formula}s को एक साथ संतुष्ट करता हो। क्योंकि !!{formula}s आगमनात्मक रूप से
परिभाषित हैं और संतुष्टि भी !!{formula}s की संरचना पर आगमन द्वारा परिभाषित है,
हम अपने अर्थविज्ञान के गुणधर्म सिद्ध करने तथा परिभाषित अर्थवैज्ञानिक धारणाओं के
बीच संबंध स्थापित करने के लिए आगमन का उपयोग कर सकते हैं।


\end{document}
